\chapter{CVE-2021-3600}

\section{CVE Selection}
I selected CVE-2021-3600 inspired by previous years' research projects, such as CVE-2020-8835, which pioneered scalar bounds confusion on BPF Maps. This case provides an architectural study of how post-verification instruction rewriting can undermine the safety proofs established by the verifier, and it allowed me to test how modern kernel mitigations, specifically the \texttt{alu\_limit} pointer arithmetic sanitization, respond to a classic bounds tracking mismatch between the verifier and the JIT compiler.

\section{Vulnerability Analysis}

\subsection{Root Cause}

The root cause lies in \texttt{fixup\_bpf\_calls()} in \texttt{kernel/bpf/verifier.c}. This function is a post-verification pass that rewrites certain BPF instructions before JIT compilation. For 32-bit division and modulo operations, the vulnerable code inserts:

\textbf{Vulnerable Code and Truncation Bug}

\begin{lstlisting}[style=kernel,caption={Vulnerable fixup code (kernel/bpf/verifier.c, fixup\_bpf\_calls())}]
struct bpf_insn mask_and_div[] = {
    BPF_MOV32_REG(insn->src_reg, insn->src_reg),  // BUG!
    /* Rx div 0 -> 0 */
    BPF_JMP_IMM(BPF_JNE, insn->src_reg, 0, 2),   // 64-bit jump!
    BPF_ALU32_REG(BPF_XOR, insn->dst_reg, insn->dst_reg),
    BPF_JMP_IMM(BPF_JA, 0, 0, 1),
    *insn,
};
\end{lstlisting}

\textbf{Why \texttt{BPF\_MOV32\_REG} is wrong:}
The \texttt{BPF\_MOV32\_REG(src, src)} instruction zeroes the upper 32 bits of the source register at \emph{runtime}. However, the verifier's abstract interpretation of the program happens \emph{before} \texttt{fixup\_bpf\_calls()} runs. The verifier never sees this inserted \texttt{mov32} instruction, so its bounds tracking for the source register remains unchanged, it still reflects the full 64-bit value.

Furthermore, during the actual \texttt{BPF\_\allowbreak ALU | BPF\_\allowbreak DIV | BPF\_\allowbreak X} verification, the function \texttt{adjust\_\allowbreak scalar\_\allowbreak min\_\allowbreak max\_\allowbreak vals()} only marks the \emph{destination} register as unknown after division, the source register's bounds are not updated.

\textbf{The double bug:}
There is actually a second problem: the zero-check \texttt{BPF\_JMP\_IMM(BPF\_JNE, src, 0, 2)} uses a \textbf{64-bit comparison}, even for 32-bit operations. After the \texttt{mov32} truncation, a value like \texttt{0xFFFFFFFF} (which is nonzero) would not trigger the zero-check. But the verifier, which still tracks the original 64-bit value, may compute the wrong result for the division.

\textbf{Concrete Example:}

\begin{lstlisting}[style=kernel,caption={Bounds mismatch demonstration}]
// Source level:
0: (b7) r0 = -1           // R0 = 0xFFFFFFFF_FFFFFFFF
1: (b7) r1 = -1           // R1 = 0xFFFFFFFF_FFFFFFFF
2: (3c) w0 /= w1          // 32-bit division

// After fixup_bpf_calls() rewrites instruction 2:
2a: (bc) w1 = w1          // mov32: R1 -> 0x0000_0000_FFFF_FFFF
2b: (15) if r1 == 0 goto  // 64-bit zero check (R1 != 0, so not taken)
2c: (3c) w0 /= w1         // actual division
\end{lstlisting}

\begin{table}[H]
\centering
\begin{tabular}{@{}lll@{}}
\toprule
\textbf{Register} & \textbf{Verifier State} & \textbf{Runtime Value} \\
\midrule
R1 before div & \texttt{0xFFFFFFFF\_FFFFFFFF} (-1) & \texttt{0x0000\_0000\_FFFF\_FFFF} (after mov32) \\
R1 >> 32 & \texttt{0xFFFFFFFF} & \texttt{0x00000000} \\
\bottomrule
\end{tabular}
\caption{Verifier vs. runtime state after the truncation bug}
\end{table}

\subsection{Patch Details}

The vulnerability is addressed by correcting the instruction rewriting logic in \texttt{fixup\_bpf\_calls()}. Specifically, the patch:
\begin{enumerate}
    \item \textbf{Removes} the \texttt{BPF\_MOV32\_REG(src, src)} instruction insertion, eliminating the truncation of the source register at runtime.
    \item \textbf{Replaces} the 64-bit \texttt{JMP\_IMM} zero-check with a proper \texttt{JMP32} instruction for 32-bit operations.
\end{enumerate}

Without the inserted \texttt{mov32}, the verifier's tracked bounds always remain synchronized with the actual runtime values.

\begin{table}[H]
\centering
\begin{tabular}{@{}ll@{}}
\toprule
\textbf{Kernel series} & \textbf{Fixed in} \\
\midrule
5.4.x & 5.4.101 \\
5.10.x & 5.10.19 \\
5.11.x & 5.11.2 \\
\bottomrule
\end{tabular}
\caption{Backport versions for CVE-2021-3600}
\end{table}

\subsection{Exploiting the Bounds Mismatch}

The key insight is creating a register whose verifier-tracked value differs from its runtime value after a 32-bit division:

\begin{itemize}
    \item Set source register \texttt{R3 = 0xFFFFFFFF\_00000001}
    \item Set target register \texttt{R2 = target\_offset} (any value)
    \item Execute \texttt{w2 /= w3} (32-bit div)
    \item \textbf{Verifier}: \texttt{R2 = target / 0xFFFFFFFF\_00000001 $\approx$ 0} (very large divisor $\rightarrow$ quotient near zero)
    \item \textbf{Runtime}: \texttt{R3} truncated to \texttt{0x00000001}, so \texttt{R2 = target / 1 = target} (unchanged!)
\end{itemize}

Now the verifier believes \texttt{R2 $\approx$ 0} (within map bounds), but at runtime \texttt{R2 = target\_offset} which can be out of bounds. Using \texttt{R2} as an index into a map value pointer gives us an \textbf{out-of-bounds read or write}.


% ═══════════════════════════════════════════════════════
\section{Environment Setup}

\begin{table}[H]
\centering
\begin{tabular}{@{}ll@{}}
\toprule
\textbf{Component} & \textbf{Details} \\
\midrule
Host OS & Linux (Ubuntu-based) \\
Virtualization & QEMU/KVM (Buildroot-based rootfs) \\
Kernel Version & 5.10.15 (vulnerable) \\
Rootfs & Buildroot minimal with custom exploit overlay \\
Kernel Config & \texttt{CONFIG\_BPF\_SYSCALL=y}, \texttt{CONFIG\_BPF\_JIT=y} \\
& \texttt{CONFIG\_KALLSYMS\_ALL=y} \\
& KASLR disabled (\texttt{nokaslr}), SMEP/SMAP disabled \\
VM Resources & 2 vCPUs, 512 MB RAM \\
Compilation & GCC, static linking, \texttt{-O2} \\
User & root (for Milestones 1--3), user uid=1000 (for M4 LPE) \\
\bottomrule
\end{tabular}
\caption{Test environment for CVE-2021-3600}
\end{table}

Kernel version 5.10.15 was selected because the fix for CVE-2021-3600 was backported to the 5.10 stable branch starting from version 5.10.16. The vulnerable \texttt{mask\_\allowbreak and\_\allowbreak div[]} code with the buggy \texttt{BPF\_\allowbreak MOV32\_\allowbreak REG\allowbreak (src, \allowbreak src)} is present in 5.10.15. A matching kernel configuration file (\texttt{v5.10.15.config}) was created to ensure BPF subsystem options were correctly enabled.


% ═══════════════════════════════════════════════════════
\section{Proof of Concept (PoC)}

\subsection{Milestone 1: Demonstrate Bounds Confusion}

\textbf{Objective}
Prove that the BPF verifier's bounds tracking becomes stale after a 32-bit division due to the inserted \texttt{BPF\_MOV32\_REG(src, src)}, demonstrating that the runtime value of the source register differs from what the verifier tracks.

\textbf{Approach}
The BPF program:
\begin{enumerate}
    \item Sets \texttt{R0 = R1 = -1} (\texttt{0xFFFFFFFF\_FFFFFFFF})
    \item Saves a copy of \texttt{R1} to \texttt{R8} (before the fixup affects R1)
    \item Executes \texttt{w0 /= w1} (32-bit division)
    \item After division, right-shifts \texttt{R1 >> 32} to expose the truncation
    \item Exfiltrates the shifted value and the original \texttt{R8 >> 32} to a BPF map
\end{enumerate}

The key comparison:
\begin{itemize}
    \item \texttt{R1 >> 32} should be \texttt{0} at runtime (upper bits zeroed by mov32) but the verifier thinks it's \texttt{0xFFFFFFFF}
    \item \texttt{R8 >> 32} should be \texttt{0xFFFFFFFF} (R8 was copied before the fixup)
\end{itemize}

\textbf{Results}

\begin{lstlisting}[style=log,caption={Milestone 1 output on vulnerable kernel 5.10.15}]
[+] BPF program LOADED (fd=5) - verifier accepted!

  result[0] (div result w0/w1)   = 0x0000000000000001  <- 0xFFFFFFFF / 0xFFFFFFFF = 1
  result[1] (R1 >> 32 after div) = 0x0000000000000000  <- upper bits ZEROED (BUG!)
  result[2] (orig R1 >> 32)      = 0x00000000ffffffff  <- original had upper bits

[+] VULNERABILITY CONFIRMED:
    The verifier's fixup_bpf_calls() inserted mov32 that
    truncated R1's upper 32 bits at runtime, but the verifier
    still tracks the original 64-bit bounds.
\end{lstlisting}

\textbf{Analysis:} Result[1] being 0 while result[2] is \texttt{0xFFFFFFFF} proves the \texttt{BPF\_MOV32\_REG(R1, R1)} zeroed R1's upper bits at runtime while the verifier continued tracking the original value. The verifier's bounds are stale. This milestone confirms the vulnerability is present and exploitable on kernel 5.10.15.


% ═══════════════════════════════════════════════════════
\subsection{Milestone 2: OOB Read (KASLR Bypass)}

\textbf{Objective}
Leverage the bounds confusion to perform out-of-bounds reads past a BPF map value boundary, leaking kernel heap data.

\textbf{Approach}

\begin{enumerate}
    \item Create a BPF array map with \texttt{value\_size = 256} bytes
    \item Lookup the map value to get \texttt{PTR\_TO\_MAP\_VALUE} in \texttt{R6}
    \item Set \texttt{R2 = oob\_offset} (target offset beyond the 256-byte boundary)
    \item Set \texttt{R3 = 0xFFFFFFFF\_00000001} via 64-bit immediate load
    \item Execute \texttt{w2 /= w3}:
    \begin{itemize}
        \item Verifier: \texttt{R2 = oob\_offset / 0xFFFFFFFF\_00000001 = 0} (within bounds)
        \item Runtime: \texttt{R3} truncated to 1, so \texttt{R2 = oob\_offset} (OOB!)
    \end{itemize}
    \item Compute \texttt{R7 = R6 + R2} (map pointer + confused index)
    \item Read \texttt{*(R7 + 0)}, verifier allows (thinks offset = 0), runtime reads OOB
    \item Exfiltrate leaked data to result map
\end{enumerate}

For each of 15 OOB offsets (256, 264, 272, ..., 368), a separate program is loaded and run.

\textbf{Results}

\begin{lstlisting}[style=log,caption={Milestone 2 output on kernel 5.10.15 (as root)}]
[*] === CVE-2021-3600 Milestone 2: OOB Read via Bounds Confusion ===
[+] Data map created (fd=3, value_size=256)
[+] Result map created (fd=4)
[*] Attempting OOB reads beyond map value (boundary at +0x100)...

  leaked[ 0] (+0x100) = 0x0000000000000000
  leaked[ 1] (+0x108) = 0x0000000000000000
  ...
  leaked[14] (+0x170) = 0x0000000000000000

  Total OOB reads: 15/15 successful
\end{lstlisting}

\textbf{Analysis:} The OOB reads return \texttt{0x0000000000000000}, zero-filled slab padding beyond the 256-byte map value. This is \textbf{different} from the map fill value (\texttt{0x4141414141414141}), confirming that the BPF program genuinely accesses memory past the map value boundary. The leaked data corresponds to the SLUB allocator's zero-initialized padding after the map value allocation, which in a more densely populated heap would contain kernel pointers and metadata from adjacent allocations.


% ═══════════════════════════════════════════════════════
\subsection{Milestone 3: OOB Write + Escalation Analysis}

\textbf{Objective}
Extend the bounds confusion to write operations and analyze the feasibility of privilege escalation.

\textbf{Approach}
The same bounds confusion used for reads also works for writes:
\begin{enumerate}
    \item Obtain the confused index \texttt{R2} (verifier: 0, runtime: OOB offset)
    \item Compute \texttt{R7 = map\_ptr + R2}
    \item Write a marker value: \texttt{*(R7) = 0xDEAD1337CAFE4242}
    \item Read back the written value to confirm the OOB write
    \item Run a separate program to independently verify persistence
    \item Repeat across 5 rounds for reliability testing
\end{enumerate}

\textbf{Results}
\begin{lstlisting}[style=log,caption={Milestone 3 output on kernel 5.10.15}]
Part 1: OOB Write
  Write program read-back: 0xdead1337cafe4242  <- marker confirmed
  Independent read:        0xdead1337cafe4242  <- OOB WRITE CONFIRMED

Part 2: Reliability test (5 rounds)
  Run 1/5: PASS
  Run 2/5: PASS
  ...
  Reliability: 5/5 successful (100%)
\end{lstlisting}

\textbf{Escalation Analysis}

The combined OOB read (Milestone 2) and OOB write (Milestone 3) primitives provide:

\begin{table}[H]
\centering
\begin{tabular}{@{}lll@{}}
\toprule
\textbf{Escalation Path} & \textbf{Feasibility} & \textbf{Notes} \\
\midrule
Corrupt adjacent map ops & Medium & Needs KASLR bypass (from M2) \\
Corrupt slab metadata & Hard & Heap layout dependent \\
Overwrite cred structure & Hard & Requires heap spray \\
Hijack function pointer & Medium & Chain M2 leak + M3 write \\
\bottomrule
\end{tabular}
\caption{Privilege escalation feasibility analysis}
\end{table}

The most viable path combines Milestone 2 (leaking kernel pointers to defeat KASLR) with Milestone 3 (OOB write to corrupt a function pointer table), allowing for execution redirection to \texttt{commit\_\allowbreak creds (prepare\_\allowbreak kernel\_\allowbreak cred(0))} for root.


% ═══════════════════════════════════════════════════════
\subsection{Milestone 4: Full LPE Attempt}

\textbf{Objective}
Achieve local privilege escalation from an unprivileged user (\texttt{uid=1000}) to root by exploiting the OOB read/write primitives established in Milestones 2 and 3.

\textbf{Exploit Strategy}
The designed exploit followed a well-established eBPF exploitation pattern:

\begin{enumerate}
    \item \textbf{Heap grooming}: Allocate many BPF array map pairs in rapid succession to achieve adjacent placement in the SLUB allocator
    \item \textbf{OOB read scan}: Use the first map of each pair to read beyond its boundary, searching for the \texttt{array\_map\_ops} function pointer table of the second map. This is a known signature at a fixed kernel address (\texttt{0xffffffff820201e0} on 5.10.15)
    \item \textbf{OOB write corruption}: Once adjacent maps are found, overwrite the \texttt{ops} pointer of the victim map with a userspace address pointing to a crafted fake \texttt{bpf\_map\_ops} structure
    \item \textbf{ret2usr}: The fake ops table redirects \texttt{map\_lookup\_elem} to a userspace function that calls \texttt{commit\_creds(prepare\_kernel\_cred(0))}, achieving privilege escalation
\end{enumerate}

\textbf{Implementation Details}

The exploit (\texttt{milestone4.c}) implemented:

\begin{itemize}
    \item Dynamic symbol resolution from \texttt{/proc/kallsyms} with fallback to hardcoded addresses
    \item \texttt{RLIMIT\_MEMLOCK} elevation to allow sufficient BPF map allocations
    \item Incrementing fill patterns per map pair to distinguish OOB data from in-bounds reads
    \item 16 map pairs with 256-byte values to maximize heap grooming success
    \item Diagnostic output comparing OOB-read values against the known fill pattern
\end{itemize}

\textbf{Results: Failure and Root Cause Analysis}

The OOB primitives from Milestones 2 and 3 were confirmed to work \textbf{as root}. However, executing the exploit as an unprivileged user (\texttt{uid=1000}) revealed that all OOB reads returned the map's own fill pattern, indicating no actual out-of-bounds access:

\begin{lstlisting}[style=log,caption={OOB comparison: root vs. unprivileged on kernel 5.10.15}]
# As root:
  leaked[ 0] (+0x100) = 0x0000000000000000  <- genuine OOB (zeros)
  leaked[ 1] (+0x108) = 0x0000000000000000  <- genuine OOB (zeros)

# As uid=1000 (unprivileged):
  leaked[ 0] (+0x100) = 0x4141414141414141  <- offset-0 data (NOT OOB)
  leaked[ 1] (+0x108) = 0x4141414141414141  <- offset-0 data (NOT OOB)
\end{lstlisting}

\textbf{Root Cause and Exploit Mitigations}

\textbf{1. \texttt{alu\_limit} pointer arithmetic sanitization:}

For unprivileged BPF programs, the verifier applies additional pointer arithmetic sanitization via the \texttt{alu\_limit} mechanism (in \texttt{fixup\_bpf\_calls()}, lines 10921--10970 of \texttt{verifier.c}). When an \texttt{ALU64\_ADD} instruction adds a scalar to a pointer (\texttt{PTR + SCALAR}), the verifier inserts bounds-clamping code that constrains the scalar offset within the verifier's tracked bounds. Since the verifier believes the confused index is $\approx 0$, the runtime sanitization clamps the actual OOB offset back to 0, neutralizing the exploit.

This sanitization is \emph{only} applied to unprivileged programs. Privileged programs (\texttt{CAP\_BPF}/root) bypass the \texttt{alu\_limit} checks entirely (via \texttt{env->bypass\_spec\_v1 = true}), which is why the OOB confusion works as root but not as an unprivileged user.

\textbf{2. Double-division pattern:}

I also attempted a \emph{double-division} variant, hypothesizing that consecutive sequences of \texttt{LD\_IMM64}, \texttt{MOV32\_IMM}, and \texttt{ALU32\_DIV} would bypass the verifier's state pruning and the \texttt{alu\_limit} sanitization. This approach did not improve results because each division independently triggers \texttt{BPF\_MOV32\_REG} truncation, and the \texttt{alu\_limit} sanitization applies independently to the final pointer addition regardless of how the scalar was computed.

\textbf{3. \texttt{/proc/kallsyms} access barrier:}

The exploit attempted to dynamically resolve kernel symbols (including \texttt{commit\_\allowbreak creds}, \texttt{prepare\_\allowbreak kernel\_\allowbreak cred}, and \texttt{array\_\allowbreak map\_\allowbreak ops}) from \texttt{/proc/kallsyms}. However, unprivileged users receive zeroed addresses due to the \texttt{kptr\_restrict} sysctl setting. This was mitigated by hardcoding addresses extracted as root, but represents yet another barrier to fully unprivileged exploitation.


% ═══════════════════════════════════════════════════════
\section{Conclusion}

\subsection{Exploitation Impact Analysis}

\begin{table}[H]
\centering
\begin{tabular}{@{}llcc@{}}
\toprule
\textbf{Kernel} & \textbf{Privilege} & \textbf{OOB Works?} & \textbf{LPE Possible?} \\
\midrule
5.10.15 & root & \ding{51} & N/A (already root) \\
5.10.15 & unprivileged & \ding{55} (alu\_limit) & \ding{55} \\
\bottomrule
\end{tabular}
\caption{Exploitability matrix for CVE-2021-3600}
\end{table}



CVE-2021-3600 requires \texttt{CAP\_BPF} (or root) to be exploited because:

\begin{enumerate}
    \item The vulnerability creates a \emph{verifier/runtime bounds mismatch} for the source register of a 32-bit division
    \item To exploit this mismatch, the confused register must be used as an index for \texttt{PTR + SCALAR} arithmetic
    \item For \textbf{unprivileged programs}, the verifier inserts \texttt{alu\_limit} sanitization on all \texttt{PTR + SCALAR} operations. This sanitization uses the verifier's tracked bounds (which believe the index is $\approx 0$) to clamp the runtime value, \emph{independently of the actual register value}
    \item For \textbf{privileged programs}, the \texttt{alu\_limit} sanitization is skipped (\texttt{env->bypass\_spec\_v1 = true}), so the confused index passes through unclamped
\end{enumerate}

This creates a paradox: the exploit only works for users who already have \texttt{CAP\_BPF}, which typically implies root-equivalent privileges. The vulnerability cannot be used for privilege escalation from an unprivileged user.




% ═══════════════════════════════════════════════════════
\subsection{Failed Approaches \& Lessons Learned}

\textbf{Failed: OOB Confusion as Unprivileged User}

The OOB confusion worked as root but not as an unprivileged user. Investigation revealed the \texttt{alu\_limit} pointer arithmetic sanitization that clamps scalar offsets for unprivileged programs. This runtime sanitization operates independently from the verifier's abstract bounds tracking, creating a second layer of defense that neutralizes the bounds confusion even when the verifier's view is successfully corrupted.

\textbf{Lesson}: The BPF verifier has multiple layers of defense. The bounds confusion bypasses the verifier's abstract interpretation, but a separate runtime sanitization mechanism (\texttt{alu\_limit}) independently enforces bounds for unprivileged users. Exploiting such vulnerabilities requires understanding all defense layers.

\textbf{Failed: Double-Division Bypass}

I hypothesized that performing two consecutive 32-bit divisions would confuse the verifier's state pruning and bypass the \texttt{alu\_limit} sanitization. The double-div pattern did not change the outcome because \texttt{alu\_limit} operates on the final \texttt{PTR + SCALAR} operation, not on the scalar computation itself.

\textbf{Lesson}: Defense mechanisms that operate at the point-of-use (pointer arithmetic sanitization) are more robust than those that rely on tracking intermediate computations.

\textbf{Failed: \texttt{/proc/kallsyms} Symbol Resolution}

The exploit attempted to dynamically resolve kernel symbols from \texttt{/proc/kallsyms}. However, unprivileged users receive zeroed addresses due to \texttt{kptr\_restrict} settings. This was mitigated by hardcoding addresses extracted as root, but demonstrates an additional barrier to unprivileged exploitation.

\textbf{Lesson}: Real-world kernel exploitation from unprivileged contexts must contend with multiple independent hardening mechanisms beyond the specific vulnerability being targeted.


% ═══════════════════════════════════════════════════════
